\documentclass[conference]{IEEEtran}
\IEEEoverridecommandlockouts

\usepackage[brazilian]{babel}
\usepackage[utf8]{inputenc}
\usepackage[T1]{fontenc}
\usepackage{cite}
\usepackage{amsmath,amssymb}
\usepackage{booktabs}
\usepackage{array}
\usepackage{url}
\usepackage{hyperref}
\usepackage{xspace}
\renewcommand{\IEEEkeywordsname}{Palavras-chave}

% Números citados no texto: gerados por paper/gera_tabelas.py a partir de sim/results/.
% Gerado por paper/gera_tabelas.py. Não editar à mão.
\newcommand{\NumMatrizes}{8\xspace}
\newcommand{\NumExecSpmv}{96\xspace}
\newcommand{\NumBitExatas}{96\xspace}
\newcommand{\MinLinhas}{30\xspace}
\newcommand{\MaxLinhas}{1033\xspace}
\newcommand{\MaxNnz}{6027\xspace}
\newcommand{\MedFPtrinta}{$5{,}91 \times 10^{-8}$\xspace}
\newcommand{\MedTakumTrinta}{$1{,}15 \times 10^{-8}$\xspace}
\newcommand{\MedFPtrintaAcSes}{$3{,}61 \times 10^{-8}$\xspace}
\newcommand{\MedTakumTrintaAcTrinta}{$2{,}94 \times 10^{-8}$\xspace}
\newcommand{\RazaoFormatoAcTrinta}{2{,}01\xspace}
\newcommand{\RazaoAcumuladorTakum}{2{,}56\xspace}
\newcommand{\RazaoKernelTrinta}{5{,}15\xspace}
\newcommand{\RazaoAcAmpliadoTrinta}{3{,}15\xspace}
\newcommand{\ParTrintaFPAmbos}{8\xspace}
\newcommand{\ParTrintaFPTakumVence}{7\xspace}
\newcommand{\ParTrintaFPOutroVence}{1\xspace}
\newcommand{\ParTrintaFPEmpates}{0\xspace}
\newcommand{\ParTrintaFPTakumNF}{0\xspace}
\newcommand{\ParTrintaFPOutroNF}{0\xspace}
\newcommand{\ParTrintaAcFPAmbos}{8\xspace}
\newcommand{\ParTrintaAcFPTakumVence}{7\xspace}
\newcommand{\ParTrintaAcFPOutroVence}{1\xspace}
\newcommand{\ParTrintaAcFPEmpates}{0\xspace}
\newcommand{\ParTrintaAcFPTakumNF}{0\xspace}
\newcommand{\ParTrintaAcFPOutroNF}{0\xspace}
\newcommand{\ParDezFPAmbos}{5\xspace}
\newcommand{\ParDezFPTakumVence}{3\xspace}
\newcommand{\ParDezFPOutroVence}{2\xspace}
\newcommand{\ParDezFPEmpates}{0\xspace}
\newcommand{\ParDezFPTakumNF}{0\xspace}
\newcommand{\ParDezFPOutroNF}{3\xspace}
\newcommand{\ParOitoEquatromAmbos}{4\xspace}
\newcommand{\ParOitoEquatromTakumVence}{0\xspace}
\newcommand{\ParOitoEquatromOutroVence}{4\xspace}
\newcommand{\ParOitoEquatromEmpates}{0\xspace}
\newcommand{\ParOitoEquatromTakumNF}{0\xspace}
\newcommand{\ParOitoEquatromOutroNF}{4\xspace}
\newcommand{\ParOitoCincodoisAmbos}{5\xspace}
\newcommand{\ParOitoCincodoisTakumVence}{3\xspace}
\newcommand{\ParOitoCincodoisOutroVence}{2\xspace}
\newcommand{\ParOitoCincodoisEmpates}{0\xspace}
\newcommand{\ParOitoCincodoisTakumNF}{0\xspace}
\newcommand{\ParOitoCincodoisOutroNF}{3\xspace}
\newcommand{\NFTakumOito}{0\xspace}
\newcommand{\NFTakumDez}{0\xspace}
\newcommand{\NFEquatrom}{4\xspace}
\newcommand{\NFCincodois}{3\xspace}
\newcommand{\NFDez}{3\xspace}
\newcommand{\ElemNnzFP}{5{,}33\xspace}
\newcommand{\ElemNnzTakum}{4{,}50\xspace}
\newcommand{\InsnNnzFPCento}{6{,}14\xspace}
\newcommand{\CsrNnzFPCento}{1{,}33\xspace}
\newcommand{\InsnNnzFPDuzentos}{4{,}59\xspace}
\newcommand{\CsrNnzFPDuzentos}{1{,}07\xspace}
\newcommand{\InsnNnzFPQuinhentos}{4{,}34\xspace}
\newcommand{\CsrNnzFPQuinhentos}{1{,}03\xspace}
\newcommand{\InsnNnzTrintaCento}{5{,}98\xspace}
\newcommand{\CsrNnzTrintaCento}{1{,}33\xspace}
\newcommand{\RelInsnTrintaCento}{$-$2{,}7\%\xspace}
\newcommand{\RelCsrTrintaCento}{+0{,}0\%\xspace}
\newcommand{\EconomiaCento}{2{,}7\%\xspace}
\newcommand{\ZvtNnzTrintaCento}{0{,}50\xspace}
\newcommand{\InsnNnzTrintaDuzentos}{4{,}55\xspace}
\newcommand{\CsrNnzTrintaDuzentos}{1{,}07\xspace}
\newcommand{\RelInsnTrintaDuzentos}{$-$0{,}7\%\xspace}
\newcommand{\RelCsrTrintaDuzentos}{+0{,}0\%\xspace}
\newcommand{\EconomiaDuzentos}{0{,}7\%\xspace}
\newcommand{\ZvtNnzTrintaDuzentos}{0{,}37\xspace}
\newcommand{\InsnNnzTrintaQuinhentos}{4{,}33\xspace}
\newcommand{\CsrNnzTrintaQuinhentos}{1{,}03\xspace}
\newcommand{\RelInsnTrintaQuinhentos}{$-$0{,}3\%\xspace}
\newcommand{\RelCsrTrintaQuinhentos}{+0{,}0\%\xspace}
\newcommand{\EconomiaQuinhentos}{0{,}3\%\xspace}
\newcommand{\ZvtNnzTrintaQuinhentos}{0{,}35\xspace}
\newcommand{\InsnNnzDezCento}{4{,}95\xspace}
\newcommand{\CsrNnzDezCento}{1{,}47\xspace}
\newcommand{\RelInsnDezCento}{$-$19{,}3\%\xspace}
\newcommand{\RelCsrDezCento}{+10{,}7\%\xspace}
\newcommand{\ZvtNnzDezCento}{0{,}37\xspace}
\newcommand{\InsnNnzDezDuzentos}{4{,}69\xspace}
\newcommand{\CsrNnzDezDuzentos}{1{,}39\xspace}
\newcommand{\RelInsnDezDuzentos}{+2{,}2\%\xspace}
\newcommand{\RelCsrDezDuzentos}{+29{,}9\%\xspace}
\newcommand{\ZvtNnzDezDuzentos}{0{,}35\xspace}
\newcommand{\InsnNnzDezQuinhentos}{4{,}60\xspace}
\newcommand{\CsrNnzDezQuinhentos}{1{,}36\xspace}
\newcommand{\RelInsnDezQuinhentos}{+6{,}0\%\xspace}
\newcommand{\RelCsrDezQuinhentos}{+32{,}5\%\xspace}
\newcommand{\ZvtNnzDezQuinhentos}{0{,}34\xspace}
\newcommand{\InsnNnzOitoCento}{4{,}51\xspace}
\newcommand{\CsrNnzOitoCento}{1{,}39\xspace}
\newcommand{\RelInsnOitoCento}{$-$26{,}6\%\xspace}
\newcommand{\RelCsrOitoCento}{+4{,}6\%\xspace}
\newcommand{\ZvtNnzOitoCento}{0{,}35\xspace}
\newcommand{\InsnNnzOitoDuzentos}{4{,}43\xspace}
\newcommand{\CsrNnzOitoDuzentos}{1{,}36\xspace}
\newcommand{\RelInsnOitoDuzentos}{$-$3{,}4\%\xspace}
\newcommand{\RelCsrOitoDuzentos}{+27{,}4\%\xspace}
\newcommand{\ZvtNnzOitoDuzentos}{0{,}34\xspace}
\newcommand{\InsnNnzOitoQuinhentos}{4{,}39\xspace}
\newcommand{\CsrNnzOitoQuinhentos}{1{,}35\xspace}
\newcommand{\RelInsnOitoQuinhentos}{+1{,}1\%\xspace}
\newcommand{\RelCsrOitoQuinhentos}{+31{,}2\%\xspace}
\newcommand{\ZvtNnzOitoQuinhentos}{0{,}34\xspace}
\newcommand{\InsnMenorTrinta}{13\xspace}
\newcommand{\InsnIgualTrinta}{11\xspace}
\newcommand{\InsnMaiorTrinta}{0\xspace}
\newcommand{\InsnMenorDez}{11\xspace}
\newcommand{\InsnIgualDez}{0\xspace}
\newcommand{\InsnMaiorDez}{13\xspace}
\newcommand{\InsnMenorOito}{12\xspace}
\newcommand{\InsnIgualOito}{0\xspace}
\newcommand{\InsnMaiorOito}{12\xspace}
\newcommand{\InsnPares}{24\xspace}
\newcommand{\ModRFP}{13\xspace}
\newcommand{\ModCFP}{12\xspace}
\newcommand{\ModKFP}{21\xspace}
\newcommand{\ModROito}{14\xspace}
\newcommand{\ModCOito}{12\xspace}
\newcommand{\ModKOito}{21\xspace}
\newcommand{\ModRDez}{14\xspace}
\newcommand{\ModCDez}{13\xspace}
\newcommand{\ModKDez}{21\xspace}
\newcommand{\ModRTrinta}{14\xspace}
\newcommand{\ModCTrinta}{11\xspace}
\newcommand{\ModKTrinta}{21\xspace}
\newcommand{\ModExatas}{96\xspace}
\newcommand{\ModExec}{96\xspace}
\newcommand{\ModSinteticosOk}{12\xspace}
\newcommand{\ModSinteticosTotal}{12\xspace}
\newcommand{\ModResiduoMax}{0\xspace}


% TODO(equipe): preencher autores, afiliação e e-mails antes da entrega.
\author{Grupo do Projeto Integrador AOC -- UFMG\\
Departamento de Engenharia Eletrônica\\
Belo Horizonte, MG, Brasil}

\title{Aritmética Takum na Extensão Vetorial RISC-V: codificação, simulação
funcional e avaliação de SpMV}

\begin{document}
\maketitle

\begin{abstract}
Extensões vetoriais modernas acumulam formatos de ponto flutuante de baixa precisão, cada um com
conversões, casos especiais e decodificadores próprios. A aritmética Takum oferece precisão
afunilada, faixa dinâmica limitada e uma decodificação comum de 12 bits. Este trabalho propõe a
extensão experimental Zvtakum, sobre a RVV 1.0, com cinco instruções para T8, T16 e T32 codificadas
no espaço custom-0. A aritmética é implementada em inteiros, com arredondamento para o par no código.
Um simulador funcional RV64IM, Zicsr, RVV e Zvtakum, em C++17, foi verificado contra o QEMU, byte a
byte, nas instruções padrão, e contra um modelo Python de racionais exatos, nas instruções Takum e nos
kernels. Avaliamos o produto matriz-vetor esparso (SpMV) em oito matrizes Harwell--Boeing, com FP32 e
Takum, para VLEN de 128 a 512 bits. As \NumExecSpmv\ execuções dos kernels são bit a bit iguais ao
modelo de referência, e o número de instruções segue um modelo analítico que reproduz todas as
execuções sem resíduo. Com acumuladores de 32 bits nos dois casos, o erro mediano do Takum32 é
\RazaoFormatoAcTrinta\ vezes menor que o do FP32. Em nenhum dos \InsnPares\ pares matriz$\times$VLEN o
kernel Takum32 executa mais instruções que o FP32, e as escritas de CSR são iguais. Em 8 e 16 bits,
o Takum não produz resultados não finitos em nenhuma das matrizes, enquanto E4M3, E5M2 e FP16 estouram
em três ou quatro delas. Onde nenhum dos formatos estoura, o resultado varia entre matrizes e formatos.
\end{abstract}

\begin{IEEEkeywords}
RISC-V, RVV, aritmética Takum, Zvtakum, SpMV, simulador funcional, custom-0.
\end{IEEEkeywords}

\section{Introdução}
Os conjuntos de instruções vetoriais acumulam formatos de ponto flutuante de baixa precisão, como
bfloat16, FP16 e os formatos OFP8 E4M3 e E5M2, voltados a inteligência artificial e computação de alto
desempenho. Cada formato adiciona conversões, casos especiais e lógica de decodificação. A RVV 1.0 não
define FP8, e o suporte a FP16 e bfloat16 depende de extensões separadas~\cite{riscv-v-spec,ofp8,bfloat16}.

A aritmética Takum~\cite{hunhold2024takum} é um formato de precisão afunilada. Seus códigos de $n$ bits
compartilham uma decodificação de 12 bits, a faixa dinâmica é limitada e não há infinito. Hunhold
\cite{hunhold2025avx102} mostrou, para o AVX10.2, que substituir formatos de ponto flutuante por Takum
simplifica o conjunto de instruções. Este trabalho faz um estudo análogo para a RVV.

As contribuições são:
\begin{itemize}
\item uma codificação para cinco instruções Takum (\texttt{vtadd}, \texttt{vtsub}, \texttt{vtmul},
\texttt{vtdot} e \texttt{vncvt}) no espaço custom-0, verificada contra o inventário do riscv-opcodes;
\item uma aritmética Takum em inteiros, sem \texttt{double}, com arredondamento para o par no código;
\item um simulador funcional RV64IM, Zicsr, RVV e Zvtakum, verificado contra o QEMU e contra um modelo de
racionais exatos;
\item kernels SpMV em RISC-V para FP32 e Takum, com o mesmo perímetro de contagem e saída verificada
bit a bit;
\item um modelo analítico de contagem de instruções dos kernels, medido com programas sintéticos e
conferido nas execuções reais;
\item controles de acumulação que separam o efeito da largura do acumulador do efeito do formato.
\end{itemize}

O escopo é funcional. Não há síntese, área, frequência, energia nem tempo de execução. A variante
logarítmica e a FP16 vetorial não são executadas nesta versão.

\section{Fundamentação teórica}
\subsection{Takum linear}
Um código Takum de $n$ bits tem sinal $S$, direção $D$, três bits de regime $R$, $r$ bits de característica
$C$ e $p = n - 5 - r$ bits de fração $F$, com $r = R$ se $D = 1$ e $r = 7 - R$ se $D = 0$. O valor é
$(1 - 3S + F/2^p)\cdot 2^e$, com $e = (-1)^S (c + S)$. Para $n < 12$, a janela de decodificação tem 12 bits,
com zeros à direita. O código $1000\ldots0$ é NaR, e os resultados fora da faixa são saturados.

\subsection{Extensão vetorial RISC-V}
A RVV 1.0 controla o comprimento dos vetores com \texttt{vtype} (SEW, LMUL, políticas de cauda e máscara)
e \texttt{vl}. A instrução \texttt{vsetvli} escolhe \texttt{vl} como o mínimo entre o AVL e VLMAX. Com
rs1 $=$ x0 e rd $\neq$ x0, o AVL é VLMAX; com rs1 $=$ rd $=$ x0, o AVL é o \texttt{vl} atual. O gather
indexado \texttt{vluxei32} lê elementos da memória a partir de índices de 32 bits.

\subsection{SpMV em CSR}
No formato CSR, a linha $i$ da matriz ocupa as posições de \texttt{rowptr[i]} a \texttt{rowptr[i+1]}, com
índices de coluna em \texttt{colidx}. O produto $y = Ax$ acumula, linha a linha, os produtos dos valores
pelos elementos de $x$ indicados por \texttt{colidx}.

\section{Trabalhos relacionados}
Hunhold et al.~\cite{hunhold2024takum,hunhold2025avx102} definem o formato Takum e avaliam sua substituição
no AVX10.2; a descrição de um decodificador em VHDL está em~\cite{hunhold2024vhdl}. Posits~\cite{gustafson2017posit},
bfloat16~\cite{bfloat16} e OFP8~\cite{ofp8} são os formatos de referência de baixa precisão. A biblioteca
libtakum~\cite{libtakum} é a implementação de referência em C. A ferramenta MuFoLAB~\cite{mufolab}
automatiza o benchmark de representação com 1401 matrizes; este trabalho não a executou (ver a seção de
limitações). A especificação da RVV~\cite{riscv-v-spec} define o comportamento padrão usado nas verificações,
com QEMU~\cite{qemu} e LLVM~\cite{llvm-mc} como ferramentas de referência.

\section{Arquitetura proposta}
\subsection{Escopo do processador}
O núcleo modelado é RV64 com as extensões I, M e Zicsr, um subconjunto da RVV 1.0 (configuração de vetor,
cargas e armazenamentos unitários e indexados, FP32 com \texttt{vfmul} e \texttt{vfredusum}) e Zvtakum.
Não há modo privilegiado, MMU, instruções comprimidas nem cache. Os contadores medem o que o programa
executa, não o tempo.

\subsection{Codificação}
As instruções Zvtakum usam o espaço custom-0 (\texttt{opcode = 0x0B}), reservado pela ISA para extensões
não padronizadas. O campo \texttt{funct6} seleciona a operação, \texttt{vm} (bit 25) segue a convenção da
RVV e \texttt{funct3} vale zero (Tabela~\ref{tab:codificacao}).

\begin{table}[t]
\caption{Codificações Zvtakum (custom-0). Campos: funct6 em 31--26, vm em 25, vs2 em 24--20, vs1 em 19--15,
funct3 em 14--12 e vd em 11--7.}
\label{tab:codificacao}
\centering
\small
\begin{tabular}{@{}llp{4.2cm}@{}}
\toprule
funct6 & Instrução & Operação \\
\midrule
0 & \texttt{vtadd.vv} & $vd_i = vs2_i + vs1_i$ \\
1 & \texttt{vtsub.vv} & $vd_i = vs2_i - vs1_i$ \\
2 & \texttt{vtmul.vv} & $vd_i = vs2_i \times vs1_i$ \\
3 & \texttt{vtdot.vv} & $vd_0 \leftarrow vd_0 + \sum_i vs2_i \times vs1_i$ \\
4 & \texttt{vncvt.t.t} & estreitamento de $2\cdot$SEW para SEW \\
\bottomrule
\end{tabular}
\end{table}

A versão anterior do projeto usava o espaço OP-V com \texttt{funct6 = 0b111001}, sob a premissa de que seria o
único valor livre. O inventário do riscv-opcodes~\cite{riscv-opcodes} mostra que esse valor está livre, mas há
entre 11 e 33 outros valores livres por classe de operando. A premissa não se sustenta, e a codificação foi
migrada.

\subsection{Semântica}
Todas as instruções respeitam \texttt{vl}, \texttt{vstart}, a máscara \texttt{v0} e as políticas \texttt{vta} e
\texttt{vma}. Com \texttt{vl = 0}, não há elementos ativos, e os elementos de cauda seguem \texttt{vta}. Elementos agnósticos recebem 1s. SEW vem de
\texttt{vtype} (8, 16 ou 32), e SEW $=64$ é ilegal nas instruções Zvtakum. O acumulador de \texttt{vtdot.vv}
tem EEW $=2\cdot$SEW e ocupa um único registrador; só o elemento 0 é escrito, e \texttt{vstart} diferente de
zero é ilegal. A origem de \texttt{vncvt.t.t} é um grupo de EMUL $=2\cdot$LMUL alinhado. Não há modo de
arredondamento nem flags acumuladas: as instruções Takum não leem \texttt{frm} nem gravam \texttt{fflags}.

\subsection{Arredondamento}
O arredondamento é feito no código, e não no valor. Para um valor exato não nulo, o simulador monta a cadeia
infinita do código, toma o prefixo de $n+1$ bits e o bit \emph{sticky}, que indica bits não nulos depois do
prefixo. Se o último bit do prefixo é 1 e o sticky está ativo, ou se o código de $n$ bits é ímpar, o resultado
é arredondado para cima. O resultado é limitado a $[1, \mathrm{maxpos}]$, e os negativos recebem o complemento
de dois. Essa regra coincide com o modelo Python para T8 completo e para amostras de T16, T32 e T64. Ela pode
diferir do vizinho mais próximo no valor, como registrado em \texttt{docs/devlog.md}.

\subsection{Unidade Takum}
A decodificação lê a janela de 12 bits, obtém S, D, R e C, e extrai a fração por um deslocamento de largura
variável. O valor é representado por expoente, significando inteiro normalizado e bit sticky. A soma alinha o
menor operando com sticky, soma ou subtrai e normaliza. O produto multiplica os significandos em até 28 bits
por 28 bits (T32) ou 60 bits por 60 bits (T64), com aritmética de 128 bits. O \texttt{vtdot} arredonda o
produto em $2n$ e depois a soma em $2n$. O estreitamento decodifica exatamente e codifica no formato menor. O
uso de \texttt{double} foi evitado de propósito: as 128 divergências registradas no modelo Python contra a
libtakum vêm de arredondamento duplo nesse intermediário.

\section{Implementação e verificação}
\subsection{Simulador}
O simulador \texttt{zt-sim} tem três partes: a aritmética Takum (\texttt{takum.cpp}), o núcleo com escalar, CSRs,
RVV e Zvtakum (\texttt{machine.cpp}) e a linha de comando, que carrega ELF64 RISC-V e atende \texttt{write} e
\texttt{exit} como o QEMU. Os contadores registram instruções por classe, escritas de CSR (incluindo \texttt{vl}
e \texttt{vtype} escritos por \texttt{vsetvl}), elementos ativos, bytes lidos e escritos e bytes de índice.

\subsection{Níveis de verificação}
\begin{enumerate}
\item \textbf{Núcleo Takum.} Tabelas T8 completas de soma, subtração e produto; narrow completo de T16 para T8;
operações amostradas de T16, T32 e T64; \texttt{vtdot} amostrado; ida e volta do codificador. São 21
verificações, todas aprovadas, contra o modelo Python.
\item \textbf{Escalar e RVV padrão.} Dois programas: operações escalares, com 28\,856 bytes de saída, e RVV
padrão, com 342 instruções por execução. O RVV inclui uma redução com o elemento 0 mascarado, políticas de cauda
e máscara \texttt{tu/mu} e \texttt{ta/ma}, \texttt{vl = 0} e \texttt{vstart} escrito pelo programa. Os dois são
executados em VLEN de 128, 256 e 512 bits, e a saída e o código de saída são comparados com o QEMU 10.2.2 em modo
usuário, com \texttt{vext\_spec=v1.0}. O QEMU é executado com \texttt{rvv\_ta\_all\_1s} e
\texttt{rvv\_ma\_all\_1s}, para que os elementos agnósticos recebam 1s, como no simulador. A RVV 1.0 permite as
duas escolhas; sem essa opção, as diferenças aparecem apenas nos elementos agnósticos. As seis execuções são iguais.
\item \textbf{Zvtakum.} Um programa com 101 instruções, 10 delas Zvtakum: as cinco operações, com máscara \texttt{v0},
cauda, \texttt{vl} parcial, \texttt{vl = 0}, grupos LMUL$=2$ e a sequência \texttt{vtdot}/\texttt{vncvt}. A saída
esperada vem do modelo Python. As três execuções (uma por VLEN) são iguais.
\item \textbf{Kernels SpMV.} Cada um dos \NumExecSpmv\ kernels (oito matrizes, quatro formatos, três VLENs) é
comparado com \texttt{scalar\_spmv}, que arredonda cada produto e cada soma no formato de acumulação, na mesma
ordem. As \NumBitExatas\ execuções são bit a bit iguais.
\item \textbf{Modelo de contagem.} Os parâmetros do modelo da Seção~\ref{sec:contagem} são medidos com
\ModSinteticosTotal\ programas sintéticos e conferidos contra as \NumExecSpmv\ execuções reais. Todas as
\ModSinteticosOk\ verificações sintéticas passam, e o resíduo máximo nas execuções reais é \ModResiduoMax.
\end{enumerate}

\section{Metodologia experimental}
As matrizes são oito da coleção Harwell--Boeing~\cite{duff1989hb}, disponíveis em
\texttt{python/data/harwell\_boeing}: \texttt{bcsstk01}, \texttt{bcsstk02}, \texttt{jpwh\_991}, \texttt{pores\_1},
\texttt{sherman1}, \texttt{well1033}, \texttt{west0067} e \texttt{west0132}. Têm entre \MinLinhas\ e \MaxLinhas\
linhas e até \MaxNnz\ não nulos. O vetor $x$ é pseudoaleatório (semente 2026), com elementos $k/2^{52}$ e
$|k| \le 2^{52}$, para forçar arredondamento em FP32 e em takum32. Uma versão anterior, com elementos de 20 bits,
dava erro zero em FP32 e em takum32 para \texttt{jpwh\_991}: os valores da matriz são inteiros pequenos, e para
essas magnitudes qualquer múltiplo de $2^{-20}$ cabe exatamente em FP32. Esse caso foi descartado como teste fraco.

Os kernels seguem o mesmo laço escalar sobre as linhas do CSR, com faixas (\emph{strip-mining}) de VLMAX
elementos. O kernel FP32 usa \texttt{vfmul.vv} e \texttt{vfredusum.vs}, com acumulador de 32 bits. O kernel takum
usa \texttt{vtdot.vv} com acumulador de $2n$ bits e \texttt{vncvt.t.t} no fim de cada linha. Os dois usam o gather
\texttt{vluxei32} para $x$, com índices de 32 bits. Os deslocamentos de $x$, pré-multiplicados pelo tamanho do
elemento, são preparados no host e são iguais para formatos do mesmo tamanho. Todo o restante executa no simulador,
inclusive os \texttt{vsetvli} e a zeragem do acumulador.

O erro é o erro relativo L2 em relação à referência binary128 (113 bits de significando), que arredonda cada produto
e cada soma. Essa referência é um modelo de software validado contra o \texttt{\_\_float128} do GCC, com 6\,014 de
6\,017 resultados iguais em bits; as três diferenças são de sinal de zero, previstas no contrato do modelo. Os
controles de acumulação repetem a mesma matriz e o mesmo $x$, com FP32 e takum32 acumulando na largura do formato
e em largura ampliada. E4M3, E5M2 e FP16 existem apenas como modelo numérico, sem contagem de instruções.

As comparações entre formatos são pareadas por matriz: contamos em quantas matrizes cada formato tem erro menor.
Medianas servem apenas para resumir uma coluna, nunca para comparar conjuntos de matrizes diferentes. Com um único
vetor $x$ e oito matrizes, os resultados devem ser lidos como ordens de grandeza, sem intervalos de confiança.

\section{Resultados e discussão}
\subsection{Corretude}
Todas as verificações da seção anterior passaram. As \NumExecSpmv\ execuções dos kernels produziram saída bit a bit
igual ao modelo de referência.

\subsection{Precisão em 32 bits}
A Tabela~\ref{tab:precisao32} mostra o erro por matriz. A mediana do erro do kernel FP32 é \MedFPtrinta, e a do
kernel Takum32 é \MedTakumTrinta. Pareado por matriz, o Takum32 tem erro menor em \ParTrintaFPTakumVence\ de
\ParTrintaFPAmbos\ matrizes; a exceção é \texttt{pores\_1}.

A diferença entre os kernels mistura dois efeitos: o FP32 acumula em 32 bits, e o Takum32 acumula em 64 bits. A
razão entre as medianas, de \RazaoKernelTrinta, se decompõe em \RazaoFormatoAcTrinta, o efeito do formato com o
mesmo acumulador de 32 bits nos dois casos, e \RazaoAcumuladorTakum, o efeito do acumulador de 64 bits do
\texttt{vtdot}. Com acumulador de 32 bits, o Takum32 tem erro menor em \ParTrintaAcFPTakumVence\ de
\ParTrintaAcFPAmbos\ matrizes. Com acumulador de 64 bits nos dois casos, o FP32 tem mediana de \MedFPtrintaAcSes, e o
Takum32 é \RazaoAcAmpliadoTrinta\ vezes melhor. O efeito de formato é, portanto, real, mas menor que a diferença
entre os kernels.

\begin{table*}[t]
\caption{Erro relativo L2 do SpMV em 32 bits (VLEN = 256), por matriz. A coluna FP32 é o kernel vetorial (acumulador de 32 bits); as colunas de controle usam o modelo numérico com o acumulador indicado.}
\label{tab:precisao32}
\centering
\footnotesize
\begin{tabular}{@{}lrcccc@{}}
\toprule
Matriz & nnz & FP32 & Takum32 & FP32 (acum.\ 64) & Takum32 (acum.\ 32) \\
\midrule
bcsstk01 & 400 & $7{,}20 \times 10^{-8}$ & $5{,}06 \times 10^{-8}$ & $3{,}57 \times 10^{-8}$ & $6{,}22 \times 10^{-8}$ \\
bcsstk02 & 4356 & $1{,}81 \times 10^{-7}$ & $1{,}88 \times 10^{-8}$ & $4{,}63 \times 10^{-8}$ & $8{,}32 \times 10^{-8}$ \\
jpwh\_991 & 6027 & $6{,}32 \times 10^{-8}$ & $4{,}19 \times 10^{-9}$ & $3{,}47 \times 10^{-8}$ & $9{,}40 \times 10^{-9}$ \\
pores\_1 & 180 & $3{,}54 \times 10^{-8}$ & $7{,}87 \times 10^{-8}$ & $3{,}26 \times 10^{-8}$ & $5{,}44 \times 10^{-8}$ \\
sherman1 & 3750 & $5{,}82 \times 10^{-8}$ & $4{,}02 \times 10^{-9}$ & $3{,}65 \times 10^{-8}$ & $4{,}93 \times 10^{-9}$ \\
well1033 & 4732 & $5{,}99 \times 10^{-8}$ & $4{,}07 \times 10^{-9}$ & $3{,}72 \times 10^{-8}$ & $6{,}45 \times 10^{-9}$ \\
west0067 & 294 & $4{,}82 \times 10^{-8}$ & $3{,}59 \times 10^{-9}$ & $3{,}27 \times 10^{-8}$ & $4{,}84 \times 10^{-9}$ \\
west0132 & 413 & $5{,}01 \times 10^{-8}$ & $3{,}21 \times 10^{-8}$ & $4{,}00 \times 10^{-8}$ & $4{,}94 \times 10^{-8}$ \\
\midrule
Mediana & & $5{,}91 \times 10^{-8}$ & $1{,}15 \times 10^{-8}$ & $3{,}61 \times 10^{-8}$ & $2{,}94 \times 10^{-8}$ \\
\bottomrule
\end{tabular}
\end{table*}


\subsection{Precisão em 8 e 16 bits}
A Tabela~\ref{tab:precisao816} mostra os formatos de 8 e 16 bits. Takum8 e Takum16 não produzem resultados não
finitos em nenhuma das oito matrizes. Em contraste, E4M3 tem overflow em
\NFEquatrom\ matrizes, E5M2 em \NFCincodois\ e FP16 em \NFDez. Esses resultados não finitos não são números e não
entram em nenhuma comparação de erro.

Nas matrizes em que ambos os formatos são finitos, o resultado depende do par comparado:
\begin{itemize}
\item Takum16 tem erro menor que FP16 (ambos com acumulador de 32 bits) em \ParDezFPTakumVence\ de \ParDezFPAmbos\
matrizes, e FP16 é melhor nas \ParDezFPOutroVence\ restantes;
\item Takum8 tem erro menor que E5M2 (ambos com acumulador de 16 bits) em \ParOitoCincodoisTakumVence\ de
\ParOitoCincodoisAmbos\ matrizes, e E5M2 é melhor nas \ParOitoCincodoisOutroVence\ restantes;
\item E4M3 tem erro menor que Takum8 em \ParOitoEquatromOutroVence\ de \ParOitoEquatromAmbos\ matrizes.
\end{itemize}
Em 8 bits, portanto, o Takum troca precisão por ausência de overflow: nunca estoura, mas perde para E4M3 nas
matrizes em que E4M3 não estoura.

\begin{table*}[t]
\caption{Relative L2 error at 8 and 16 bits, by matrix (VLEN = 256). NF($k$) denotes $k$ non-finite rows (overflow). Accumulators: E4M3, E5M2 and Takum8 at 16 bits; FP16 and Takum16 at 32 bits.}
\label{tab:precisao816}
\centering
\footnotesize
\begin{tabular}{@{}lccccc@{}}
\toprule
Matrix & E4M3 & E5M2 & FP16 & Takum16 & Takum8 \\
\midrule
bcsstk01 & NF(48) & NF(48) & NF(48) & $3.24 \times 10^{-3}$ & $6.44 \times 10^{-1}$ \\
bcsstk02 & NF(66) & $7.31 \times 10^{-2}$ & $3.76 \times 10^{-4}$ & $1.20 \times 10^{-3}$ & $2.27 \times 10^{-1}$ \\
jpwh\_991 & $3.45 \times 10^{-2}$ & $7.58 \times 10^{-2}$ & $2.77 \times 10^{-4}$ & $3.04 \times 10^{-4}$ & $7.77 \times 10^{-2}$ \\
pores\_1 & NF(30) & NF(15) & NF(15) & $1.18 \times 10^{-3}$ & $8.58 \times 10^{-1}$ \\
sherman1 & $3.95 \times 10^{-2}$ & $1.00 \times 10^{-1}$ & $3.06 \times 10^{-4}$ & $2.28 \times 10^{-4}$ & $6.40 \times 10^{-2}$ \\
well1033 & $4.24 \times 10^{-2}$ & $8.75 \times 10^{-2}$ & $3.54 \times 10^{-4}$ & $2.62 \times 10^{-4}$ & $6.47 \times 10^{-2}$ \\
west0067 & $3.43 \times 10^{-2}$ & $8.74 \times 10^{-2}$ & $3.36 \times 10^{-4}$ & $1.76 \times 10^{-4}$ & $5.18 \times 10^{-2}$ \\
west0132 & NF(17) & NF(3) & NF(3) & $1.40 \times 10^{-3}$ & $1.21 \times 10^{0}$ \\
\bottomrule
\end{tabular}
\end{table*}


\subsection{Modelo de contagem de instruções}
\label{sec:contagem}
O número de instruções de cada kernel segue

\begin{equation}
N_f = r_f \, R + c_f \, S + K,
\end{equation}

em que $R$ é o número de linhas da matriz, $S$ é o número de faixas somado sobre as linhas não vazias e $K$ é a
preparação e o encerramento comuns. O custo $r_f$ é por linha e $c_f$ é por faixa. A Tabela~\ref{tab:modelo} mostra
os parâmetros medidos. Eles reproduzem exatamente as \ModExec\ execuções reais e as \ModSinteticosTotal\ verificações sintéticas.

\begin{table}[t]
\caption{Parâmetros do modelo de contagem $N = rR + cS + K$, em que $R$ é o número de linhas e $S$ o número de faixas das linhas não vazias. A última coluna compara o modelo com as execuções reais.}
\label{tab:modelo}
\centering
\footnotesize
\begin{tabular}{@{}lcccrr@{}}
\toprule
Format & $r$ & $c$ & $K$ & Synthetic & Exact runs \\
\midrule
FP32 & 13 & 12 & 21 & 3/3 & 24/24 \\
Takum8 & 14 & 12 & 21 & 3/3 & 24/24 \\
Takum16 & 14 & 13 & 21 & 3/3 & 24/24 \\
Takum32 & 14 & 11 & 21 & 3/3 & 24/24 \\
\bottomrule
\end{tabular}
\end{table}


A diferença do Takum32 para o FP32 sai diretamente dos parâmetros. Cada linha do Takum32 custa uma instrução a mais,
o \texttt{vncvt.t.t} de fechamento, e cada faixa custa uma a menos, pois o \texttt{vtdot} substitui \texttt{vfmul} e
\texttt{vfredusum}. Logo
\begin{equation}
N_{\mathrm{Takum32}} - N_{\mathrm{FP32}} = R - S.
\end{equation}
A identidade vale em todos os \InsnPares\ pares matriz$\times$VLEN, e os empates ocorrem exatamente quando $S = R$.
Assim, o Takum32 executa menos instruções que o FP32 em \InsnMenorTrinta\ pares, é igual em \InsnIgualTrinta, e em
nenhum par executa mais. Um caso com mais instruções é possível: basta haver linhas vazias e nenhuma linha com mais de
VLMAX elementos, porque então $R > S$.

Os kernels de 8 e 16 bits diferem do FP32 no custo por faixa. O kernel de 8 bits tem $c = 12$, igual ao FP32: o
\texttt{vtdot} economiza uma instrução por faixa, e a configuração extra do índice de 32 bits gasta outra. O de 16 bits
tem $c = 13$, porque também calcula um deslocamento a mais para os valores de 2 bytes. Com $S_b$ o número de faixas para
elementos de $b$ bits, a diferença em relação ao FP32 é $R + 12\,(S_8 - S_{32})$ para Takum8 e
$R + 13\,S_{16} - 12\,S_{32}$ para Takum16. Como $S_8 \le S_{32}$ e $S_{16} \le S_{32}$, a redução de faixas compensa
o custo de uma instrução por linha somente enquanto as linhas precisam de várias faixas. Em VLEN grande, quase todas
as linhas cabem em uma faixa, a redução desaparece e resta o custo por linha. Isso explica por que o Takum8 passa de
uma economia para um custo quando o VLEN cresce.

\subsection{Custo de instruções e escritas de CSR}
A Tabela~\ref{tab:contagens} mostra instruções, instruções Zvtakum, escritas de CSR e elementos ativos por elemento
não nulo. Em VLEN de 256 bits, o kernel FP32 executa \InsnNnzFPDuzentos\ instruções por não nulo, e o kernel
Takum32 executa \InsnNnzTrintaDuzentos, uma diferença de \RelInsnTrintaDuzentos. A diferença é de
\RelInsnTrintaCento\ em VLEN de 128 bits e de \RelInsnTrintaQuinhentos\ em VLEN de 512 bits. As escritas de CSR são
as mesmas: \CsrNnzFPDuzentos\ por não nulo para FP32 e \CsrNnzTrintaDuzentos\ para Takum32, em VLEN de 256 bits. Os
dois kernels usam um \texttt{vsetvli} por faixa e dois por linha.

Os kernels de 8 e 16 bits executam duas configurações vetoriais por faixa, e isso aumenta as escritas de CSR por não
nulo: \RelCsrOitoDuzentos\ para Takum8 e \RelCsrDezDuzentos\ para Takum16, em VLEN de 256 bits. Em instruções, a
diferença vai de \RelInsnOitoCento\ (Takum8, VLEN de 128 bits) a \RelInsnOitoQuinhentos\ (VLEN de 512 bits), e para
Takum16 de \RelInsnDezCento\ a \RelInsnDezQuinhentos. O kernel Takum32 processa menos elementos por não nulo, pois o
\texttt{vtdot} conta um elemento ativo por produto, enquanto \texttt{vfmul} e \texttt{vfredusum} contam dois. Os valores
são \ElemNnzTakum\ para Takum32 e \ElemNnzFP\ para FP32.

\begin{table*}[t]
\caption{Instructions, Zvtakum instructions, CSR writes, and active elements per nonzero, summed over the eight matrices. Active elements count each item processed by vector data instructions (load, gather, product, reduction, store).}
\label{tab:contagens}
\centering
\footnotesize
\begin{tabular}{@{}llcccc@{}}
\toprule
Format & VLEN & Instr./nnz & Zvtakum/nnz & CSR writes/nnz & Active elem./nnz \\
\midrule
FP32 & 128 & 6.140 & 0.000 & 1.328 & 5.334 \\
FP32 & 256 & 4.587 & 0.000 & 1.069 & 5.334 \\
FP32 & 512 & 4.343 & 0.000 & 1.029 & 5.334 \\
\midrule
Takum32 & 128 & 5.977 & 0.497 & 1.328 & 4.501 \\
Takum32 & 256 & 4.553 & 0.368 & 1.069 & 4.501 \\
Takum32 & 512 & 4.329 & 0.347 & 1.029 & 4.501 \\
\midrule
Takum16 & 128 & 4.954 & 0.368 & 1.470 & 4.501 \\
Takum16 & 256 & 4.690 & 0.347 & 1.389 & 4.501 \\
Takum16 & 512 & 4.605 & 0.341 & 1.363 & 4.501 \\
\midrule
Takum8 & 128 & 4.510 & 0.347 & 1.389 & 4.501 \\
Takum8 & 256 & 4.431 & 0.341 & 1.363 & 4.501 \\
Takum8 & 512 & 4.392 & 0.337 & 1.350 & 4.501 \\
\bottomrule
\end{tabular}
\end{table*}


\subsection{Discussão}
Três resultados sustentam as conclusões. Primeiro, com acumuladores de mesma largura, o Takum32 tem erro cerca de
\RazaoFormatoAcTrinta\ vezes menor que o FP32 na mediana, e vence em \ParTrintaAcFPTakumVence\ de
\ParTrintaAcFPAmbos\ matrizes. Segundo, o Takum32 não custa mais instruções que o FP32 em nenhum dos pares testados, e
custa as mesmas escritas de CSR; a economia é pequena, de até \EconomiaCento\ em VLEN de 128 bits, e vem do
\texttt{vtdot}, que une o produto e a redução, descontado o estreitamento de cada linha. Terceiro, em 8 e 16 bits, o
Takum evita overflow, mas não é consistentemente mais preciso que os formatos IEEE-like, e o kernel de 8 e 16 bits
paga uma configuração vetorial a mais por faixa.

A conclusão defensável é que o Takum32 é uma alternativa de baixo custo ao FP32 para SpMV neste subconjunto de
instruções: tem erro menor em \ParTrintaAcFPTakumVence\ de \NumMatrizes\ matrizes com acumulador de mesma largura, e
custo de instruções igual ou menor. Os ganhos de 8 e 16 bits dependem da matriz e do VLEN.

\section{Limitações}
\label{sec:limitacoes}
O trabalho tem limitações que afetam a interpretação dos resultados. As oito matrizes são pequenas e não representam
a coleção SuiteSparse, e há um único vetor $x$; os erros devem ser lidos como ordens de grandeza. O estudo de 1401
matrizes com MuFoLAB não foi executado, porque exige Julia e uma rodada longa, e os resultados de representação da
versão anterior, com 100 matrizes e referência \texttt{long double}, não foram substituídos. A variante logarítmica
não existe nesta versão, e o custo da soma logarítmica ($\Delta\pm$) não foi medido. A FP16 vetorial (Zvfh) é reconhecida pelo
montador, mas não é executada pelo simulador. FP8 não existe na RVV, e seus números são apenas numéricos. Os contadores
medem instruções e eventos, não ciclos, e não há estimativa de área ou energia. A conformidade é verificada por amostra
e por casos de borda, não por prova. A comparação com o QEMU depende de preencher os elementos agnósticos com 1s,
escolha que a especificação permite e que o simulador adota. Por fim, o kernel usa um laço simples; uma versão que evite a configuração extra do índice de 32 bits
mudaria os números de 8 e 16 bits.

Trabalhos futuros incluem executar o estudo de 1401 matrizes, implementar Zvfh e o kernel FP16 com acumulação ampliada,
decidir a variante logarítmica e medir o custo de $\Delta\pm$, reduzir as trocas de configuração do kernel de 8 e 16 bits e
comparar com uma implementação em hardware.

\section{Conclusão}
Este trabalho propôs a extensão Zvtakum sobre a RVV 1.0 para T8, T16 e T32, com codificação em custom-0, semântica
especificada e uma aritmética Takum em inteiros. O simulador funcional executa RV64IM, Zicsr, RVV e Zvtakum e foi
verificado contra o QEMU e contra um modelo de racionais exatos. Nos kernels SpMV, o Takum32 tem erro menor que o FP32
em \ParTrintaAcFPTakumVence\ de \NumMatrizes\ matrizes com acumulador de mesma largura, e nunca executa mais
instruções. O custo de instruções segue um modelo analítico, $N = rR + cS + K$, que reproduz todas as execuções. Em 8 e
16 bits, o Takum não estoura, mas não é consistentemente mais preciso que os formatos IEEE-like, e o custo depende do VLEN.

\bibliographystyle{IEEEtran}
\bibliography{refs}

\end{document}
